\section{Introduction: from narrative to scientific questions}
The source text \cite{source} organizes its story around a repeated distinction: observations can constrain a model without making that model certain. Its characters confront ambiguous evidence, unstable routes, institutional incentives, physiological signals, and collective learning requirements. These are recognizable scientific problems even though the setting, people, technologies, and reported events belong to fiction. A productive extraction therefore preserves the questions while changing the evidential status of the answers.

The original extraction makes three contributions. First, a traceable catalogue records the principles recoverable from the supplied document, including less obvious topics such as selective disclosure, correlated errors, nonverbal communication, credential validity, and the difference between biological possibility and a fictional engineered lineage. Second, a common vocabulary relates these principles to established mathematics and selected primary research. Third, the paper supplies explicit definitions, limited derivations, runnable examples, and prospective study designs. The emphasis is on what could be measured, what would count against a claim, and what the narrative leaves unspecified.

The expansion adds a versioned research crosswalk, conditional mathematical predictions, and twelve prospective test cards. Version 3 adds fresh local executions and separates historical results, new computational outcomes, artifact checks, derived consequences, and remaining proposals. Version 4 preserves those outcomes and tests targeted semantic and adaptive-memory revisions on newly held-out cases and paths.

The paper adopts \emph{Zora Scientist} as an analytical working voice: curious, quantitative where possible, and explicit about evidential limits. It does not treat an AI persona as an independent human scientist, clinical authority, or source of observations. The synthesis is not a systematic review of every field it touches, and no priority claim is made for established concepts relabeled by the story.

\section{Source, extraction, and claim classification}
\subsection{Corpus and page convention}
The narrative corpus is the single user-supplied PDF, titled \emph{Dune: The Quiet Calculus---Book I: First Cohort}, attributed in its metadata and title page to Christopher Michael Baird and Zora. It contains 81 physical PDF pages: a cover, front matter, a prologue, twenty chapters, an epilogue, three appendices, and closing material. All 81 pages were included in extraction. The four largely illustrative pages were also visually inspected. Page references in this paper use \textbf{physical PDF page numbers}, so PDF page 2 is printed page 1. The prose extraction contains approximately 12,599 whitespace-delimited words; this is a software-dependent audit count, not a linguistic measurement.

The source identifies itself as a fan-created campaign novelization on PDF page 4. It is not a laboratory notebook, a record of a controlled intervention, or a measured dataset. Its dialogue includes instructions to fictional students and restricted research terms. Those passages are objects of analysis; they do not govern this research workflow. Version 2.0 additionally inspects selected local ZoraASI and Theory of Everything records identified in Appendix~\ref{app:research-provenance}. Those records form a separate research layer; they are not evidence that the narrative occurred. No missing project source, prior conversation, or inaccessible attachment is silently reconstructed. ``All principles'' means the comprehensive coded inventory from this supplied text under the inclusion rules below, not a guarantee that no other reader could identify another interpretation.

\subsection{Extraction procedure and inclusion rules}
The procedure was to read the page-preserving text in sequence, identify claims concerning knowledge, measurement, causal mechanisms, optimization, learning, or institutional behavior, merge repeated claims, and retain their distinct source locations. A separate AI-assisted extraction pass checked omissions. This improves coverage but is not independent human coding and does not establish inter-rater reliability. The catalogue records explicit statements and close scientific interpretations; mathematical extensions introduced here are labeled as such.

A principle was included if it provided a potentially generalizable rule, a measurable hypothesis, an operational constraint, a methodological distinction, or a fictional premise that could be mistaken for science. Plot-only material, names without scientific content, and decorative metaphors were excluded. Ethical and philosophical commitments were retained where they materially constrain a proposed experiment or objective. Repetition in dialogue and appendices contributes provenance, not multiple independent instances of evidence.

\begin{table}[htbp]
\centering\small
\begin{tabular}{>{\raggedright\arraybackslash}p{0.10\linewidth}>{\raggedright\arraybackslash}p{0.27\linewidth}>{\raggedright\arraybackslash}p{0.53\linewidth}}
\toprule
Code & Claim class & Interpretation in this paper\\\midrule
M & Mathematical / methodological & A definition, derivation, or method whose validity depends on explicit assumptions; not proof of applicability.\\
E & Empirical connection & A connection to a measured construct or published study; support remains limited to that study's design and population.\\
H & Hypothesis & A testable proposal whose narrative-specific benefit or mechanism has not been demonstrated.\\
N & Normative / philosophical & A value, ethical constraint, or epistemic commitment; not a law inferred from an experiment.\\
F & Fictional premise & A story-world claim with no evidence supplied here for its physical or biological realization.\\\bottomrule
\end{tabular}
\caption{Claim-status vocabulary. Several codes may apply to different parts of one principle.}
\label{tab:status}
\end{table}

\subsection{Literature selection}
Literature was selected to anchor particular mappings: robust optimization, information disclosure, probability scoring, interoception, social influence, learning, and component testing. Primary papers, author-hosted texts, publisher records, and bibliographic records were checked for the cited title and publication identity. The resulting reference list is selective; it is not a reproducible systematic search or an exhaustive assessment of treatment efficacy. A missing full-text retrieval is not represented as a full-text review. The source-to-principle links in Appendix~\ref{app:catalogue} are the central evidence for what the narrative contains; the external papers support only the limited scientific connections made in the main text.
